% kind: home
% title: 关于本站
% nav_title: 00-关于本站
% author: 山东大学网络空间安全学院
% date: 2026-03-06 00:00:00
% summary: 一个专注学院内数学课程学习、记录与分享的空间，各门课程的笔记与阶段总结会被整理成清晰的章节，方便回顾与持续更新。

\section{导语}

欢迎来到本站！这里是一个专注学院内数学课程学习、记录与分享的空间，各门课程的笔记、练习思考与阶段总结会被整理成清晰的章节，方便各位同学回顾。本站的内容以结构化整理为主，尽量做到“短而清晰、可追溯、可扩展”。如果你是第一次来到这里，建议先从“关于本站”开始，了解整体脉络与更新节奏。之后你可以从侧边栏进入不同课程的目录，按章节阅读，或通过搜索快速定位关键词。

我希望这个网站能够成为一个长期可维护的知识库，因此在组织方式上采用了“目录一章节一内容”的结构，并在页面上做了适合阅读的排版：正文居中、视觉留白、以及清晰的层级。内容会持续更新，可能包括课程笔记、习题解析、工具使用心得、以及学习方法总结。若你发现疏漏或希望补充讨论，欢迎留言或交流建议。感谢你的来访，愿这里的内容能为你的学习之路带来一点帮助与启发。

\section{打开方式}

不同的课程在侧边栏中会有不同的编号：

\begin{itemize}
\item 必修课程对应的编号为 `AXX`
\item 限选课程对应的编号为 `BXX`
\item 针对必修的拓展内容为 `EXX`
\item 与数学工具相关的内容为 `TXX`
\end{itemize}

点击侧边栏中的选项即可展开对应课程的子目录，点击右侧区域中的具体条目即可打开正文。

\section{网站来由}

我在大一下学期刚开始时，心里就有一个想法：既然我们学院有这么多数学课程，那为什么不把所有的数学课程都整理在一起呢？这样既可以用来回顾复习，也可以不断更新迭代，一直维持下去。

最开始这个想法是完完全全基于 GitHub 的，但是我发现 GitHub 的 Markdown 语法有些不兼容，再加上很多同学可能不熟悉 GitHub 的操作，所以我就想单独建一个网站来做这件事，尽可能去简单高效地完成这件事。现在新的版本已经把内容层切换为 LaTeX 驱动，也更适合长期维护数学课程笔记。
